\begin{remark}[Fully specified rational witness at $n=12$]\label{rem:fixed-weak-witness}
Setting $n=12$ ($d=5$), $\ell=5$, $\eps=\tau=10^{-60}$, and $h_i=i/6$ ($1\le i\le 5$) in~\eqref{eq:weak-realize} specifies an explicit rational polynomial $F_\tau(x)=10^{-60}x^{12}+10^{-300}\sum_{k=0}^5(q_k/c_k)x^{2k+1}$ of exact degree twelve whose odd perturbation vanishes through fourth order in $\tau$. Direct rational evaluation (verified in \texttt{checks/\allowbreak verify\_\allowbreak weak\_\allowbreak return.py}) gives $\sigma_{F_\tau}(1)<\delta_{12}$, where $1/\delta_{12}=\allowbreak 17{,}704{,}107{,}497{,}153{,}401{,}118{,}\allowbreak 612{,}992{,}587{,}674{,}419{,}200$, and $\sup_{|x|\le 2}|F_\tau'(x)|\le\sum_{j=1}^{12}j|a_j|2^{j-1}<\kappa_{12}$. At the ten rational endpoints $h_{i,\pm}=(4i\pm 1)/24$ ($1\le i\le 5$), the Cauchy remainder~\eqref{eq:weak-cauchy-err} is strictly below $4\cdot 10^{-49}$ (and an independent real variational enclosure along $s\in[0,1]$ in the clockwise polar equation of the main paper, Appendix~A is below $2\cdot 10^{-50}$), whereas $6|q(h_{i,\pm}^2)|>3\cdot 10^{-4}$; since $6<2\pi<8$, $(\Pi_{F_\tau}(h)-h)/(h\tau^5)$ has certified endpoint signs $(+,-),(-,+),(+,-),(-,+),(+,-)$. Thus the five disjoint brackets in $(0,1]$ each contain a return zero. They lie in the weak domain since $\sigma_{F_\tau}$ is nondecreasing and $\sigma_{F_\tau}(1)<\delta_{12}$. The odd part is nonzero, so Theorem~A.1(ii) of~\cite{MainArticle} bounds the total multiplicity in that domain by five. Each zero is therefore simple. These five hyperbolic cycles (three attracting and two repelling) account for all cycles with upper heights in $\sigma_{F_\tau}(h)\le\delta_{12}$.

\end{remark}
